\subsubsection{制約に基づく設計}

\term{制約に基づく設計}{Constraint-Based Design} は、制約を使って設計空間を狭める方法である。設計空間とは、可能な設計案全体の集合である。制約は、不可能または受け入れられない案を除外する役割を持つ。

制約には、ハードウェア、ソフトウェア、データ、操作手順、インタフェース、法令、性能、費用、納期などが含まれる。制約を明確にすれば、設計探索が速くなる。一方、制約が多すぎる場合や互いに矛盾する場合は、解が見つからないこともある。
